Geometri tropikal merupakan cabang geometri aljabar di atas aljabar semiring min-plus (atau max-plus), di mana polinomial berubah menjadi fungsi \textit{piecewise-affine} dan varietas aljabar menjadi \textit{tropical hypersurfaces}. Di sisi lain, origami matematis dilandasi oleh aksioma kombinatorik serta regulasi lipatan datar melalui Teorema Kawasaki dan Teorema Maekawa. Penelitian ini bertujuan untuk memodelkan pelipatan kaku (\textit{rigid folding}) origami menggunakan isometri tropikal dan transformasi fungsi \textit{piecewise-affine}, beralih dari pendekatan geometri analitik tradisional yang menggunakan rotasi $SO(3)$ non-linear menuju ranah \textit{piecewise-affine} tropikal yang sepenuhnya linear. Pendekatan ini merumuskan transformasi keseimbangan verteks (\textit{balancing condition}) untuk membuka cara komputasional yang efisien bagi simulasi pelipatan material meta dan struktur \textit{deployable}. Isometri rotasi ortogonal akan dirumuskan ke dalam pembobotan polinomial lintasan \textit{adjacency graph} pada ruang modul min-plus berdasarkan metrik pemetaan Abel-Jacobi dan fondasi geometri proyektif tropikal.

\katakunci{Geometri Tropikal, Isometri, Lipatan Kaku, Origami Matematis, \textit{Piecewise-Affine}}
